\section{Instruction Set Architecture}

This section documents every opcode supported by the VM. Instructions are grouped into five categories matching the source layout: stack manipulation, arithmetic, logic and comparison, control flow, and registers. Each entry records the mnemonic, its stack effect using standard notation, a description of the operation, and any runtime error conditions it can trigger.

\subsection*{Notation}

Stack effects are written in the form \texttt{before~--~after}, where the rightmost element is the top of the stack. For example, \texttt{a~b~--~(a+b)} means two values are consumed and one is produced. The symbol \texttt{--} separates the state before the instruction executes from the state after. An empty left-hand side means the instruction pushes without consuming; an empty right-hand side means it consumes without pushing.

Opcodes that require a literal argument encoded in the bytecode stream carry it as an immediate word following the opcode. These are marked with \texttt{<n>} (an integer literal) or \texttt{<addr>} (a bytecode word index) in the mnemonic column.

\subsection{Stack Manipulation}

These instructions operate directly on the shape of the operand stack. They move, duplicate, or discard values without performing any computation on them.

\begin{center}
\begin{tabular}{lll}
\toprule
\textbf{Mnemonic} & \textbf{Stack Effect} & \textbf{Description} \\
\midrule
\texttt{PUSH <n>} & \texttt{-- n}       & Push the integer literal \texttt{n} onto the stack. \\
\texttt{POP}      & \texttt{a --}        & Discard the top value. \\
\texttt{DUP}      & \texttt{a -- a a}    & Duplicate the top value; both copies remain. \\
\texttt{SWAP}     & \texttt{a b -- b a}  & Swap the top two values. \\
\texttt{PEEK}     & \texttt{a -- a}      & Print the top value without consuming it. \\
\bottomrule
\end{tabular}
\end{center}

\texttt{PUSH} is the only way to introduce a literal integer into the program. Its operand is encoded as the word immediately following the opcode in the bytecode stream, so \texttt{PUSH 42} assembles to the two-word sequence \texttt{[PUSH, 42]}.

\texttt{DUP} is used when a value must both be tested (consumed by a comparison or jump) and retained for the next iteration of a loop. A typical countdown idiom reads:

\begin{lstlisting}
DUP       ; keep a copy for the jump test
JZ  <end> ; consume the copy -- exit if zero
...
JMP <top> ; loop back to DUP
\end{lstlisting}

\texttt{SWAP} is necessary when the order of operands matters for a non-commutative operation. Because \texttt{b} is popped before \texttt{a}, the code sequence \texttt{PUSH 10 / PUSH 3 / SUB} computes $10 - 3$, not $3 - 10$. \texttt{SWAP} lets a program correct the order if operands were pushed in the wrong sequence.

\texttt{PEEK} is a diagnostic aid. It prints the top of the stack to standard output but leaves the stack unchanged, making it useful for observing intermediate values during development without disturbing the program's data.

\textbf{Error conditions:} \texttt{POP}, \texttt{DUP}, and \texttt{PEEK} halt with \textit{stack underflow} if the stack is empty. \texttt{SWAP} halts with \textit{stack underflow} if fewer than two values are present. \texttt{PUSH} halts with \textit{stack overflow} if the stack already holds 256 values.

\subsection{Arithmetic}

Arithmetic instructions each consume two values from the top of the stack and push a single result. The value pushed second is treated as the left-hand operand and the value pushed first (i.e., the top) as the right-hand operand, consistent with the natural push order.

\begin{center}
\begin{tabular}{lll}
\toprule
\textbf{Mnemonic} & \textbf{Stack Effect}     & \textbf{Description} \\
\midrule
\texttt{ADD}  & \texttt{a b -- (a+b)} & Integer addition. \\
\texttt{SUB}  & \texttt{a b -- (a-b)} & Integer subtraction. \\
\texttt{MULT} & \texttt{a b -- (a*b)} & Integer multiplication. \\
\texttt{DIV}  & \texttt{a b -- (a/b)} & Integer division (truncates toward zero). \\
\texttt{MOD}  & \texttt{a b -- (a\%b)} & Remainder after integer division. \\
\bottomrule
\end{tabular}
\end{center}

All five instructions operate on signed 32-bit integers (the C \texttt{int} type on all target platforms). Overflow wraps silently under standard C integer arithmetic rules; the VM does not detect it.

\texttt{DIV} performs truncating integer division: $7 \div 2 = 3$, and $-7 \div 2 = -3$. The result is always rounded toward zero, matching C's built-in \texttt{/} operator for signed integers.

\texttt{MOD} computes the remainder consistent with \texttt{DIV}: the sign of the result matches the sign of the dividend, so $-7~\%~2 = -1$.

\textbf{Error conditions:} All five instructions halt with \textit{stack underflow} if fewer than two values are present. \texttt{DIV} additionally halts with \textit{division by zero} if the divisor (top of stack) is zero before the pop.

\subsection{Logic and Comparison}

Comparison instructions treat any non-zero integer as \textit{true} and zero as \textit{false}. They consume their operands and push either \texttt{1} (true) or \texttt{0} (false), allowing their results to be used directly as conditions for the conditional jump instructions \texttt{JZ} and \texttt{JNZ}.

\begin{center}
\begin{tabular}{lll}
\toprule
\textbf{Mnemonic} & \textbf{Stack Effect}        & \textbf{Description} \\
\midrule
\texttt{EQ}  & \texttt{a b -- (a==b)} & Push 1 if \texttt{a} equals \texttt{b}, else 0. \\
\texttt{LT}  & \texttt{a b -- (a<b)}  & Push 1 if \texttt{a} is less than \texttt{b}, else 0. \\
\texttt{GT}  & \texttt{a b -- (a>b)}  & Push 1 if \texttt{a} is greater than \texttt{b}, else 0. \\
\texttt{AND} & \texttt{a b -- (a\&\&b)} & Push 1 if both \texttt{a} and \texttt{b} are non-zero, else 0. \\
\texttt{OR}  & \texttt{a b -- (a||b)} & Push 1 if either \texttt{a} or \texttt{b} is non-zero, else 0. \\
\texttt{NOT} & \texttt{a -- (!a)}     & Push 1 if \texttt{a} is zero; push 0 otherwise. \\
\bottomrule
\end{tabular}
\end{center}

\texttt{EQ}, \texttt{LT}, and \texttt{GT} provide the full set of binary comparisons needed to express any relational condition. Less-than-or-equal ($a \leq b$) can be expressed as \texttt{GT NOT} (negate the greater-than result); greater-than-or-equal as \texttt{LT NOT}; not-equal as \texttt{EQ NOT}.

\texttt{AND} and \texttt{OR} implement \textit{logical} (not bitwise) conjunction and disjunction. Any non-zero value is treated as true, so \texttt{AND} of \texttt{5} and \texttt{3} yields \texttt{1}, not \texttt{1} (bitwise). This matches the semantics of C's \texttt{\&\&} and \texttt{||} operators.

A typical guarded conditional pattern — execute a block only if two conditions both hold — looks like:

\begin{lstlisting}
; stack: ... a b
AND          ; stack: ... (a && b)
JZ  <skip>   ; jump past the block if false
... ; body
<skip>:
\end{lstlisting}

\textbf{Error conditions:} All binary instructions (\texttt{EQ}, \texttt{LT}, \texttt{GT}, \texttt{AND}, \texttt{OR}) halt with \textit{stack underflow} if fewer than two values are present. \texttt{NOT} halts with \textit{stack underflow} if the stack is empty.

\subsection{Control Flow}

Control flow instructions alter the instruction pointer (\texttt{ip}), enabling loops, conditionals, and function calls. All jump targets are expressed as absolute word indices into the bytecode array — the same indices reported by the \texttt{LIST} shell command.

\begin{center}
\begin{tabular}{lll}
\toprule
\textbf{Mnemonic}   & \textbf{Stack Effect} & \textbf{Description} \\
\midrule
\texttt{JMP <addr>}  & \texttt{--}      & Unconditional jump to word index \texttt{addr}. \\
\texttt{JZ  <addr>}  & \texttt{a --}    & Pop \texttt{a}; jump to \texttt{addr} if \texttt{a} is zero. \\
\texttt{JNZ <addr>}  & \texttt{a --}    & Pop \texttt{a}; jump to \texttt{addr} if \texttt{a} is non-zero. \\
\texttt{CALL <addr>} & \texttt{--}      & Push return address to call stack; jump to \texttt{addr}. \\
\texttt{RET}         & \texttt{--}      & Pop return address from call stack; resume there. \\
\bottomrule
\end{tabular}
\end{center}

\texttt{JMP} sets \texttt{ip} to the given word index unconditionally. Because the assembler emits each instruction's opcode at a known word offset (reported by \texttt{[@ word N]} in the REPL), jump targets can be calculated by inspection. Note that \texttt{JMP} itself occupies two words (opcode + operand), so a backward jump to the instruction at word 4 requires the operand \texttt{4}, not an offset.

\texttt{JZ} and \texttt{JNZ} are the conditional branches. They always consume the top of the stack regardless of whether the jump is taken. This means the condition value must be produced fresh on each iteration of a loop — typically via \texttt{DUP} before a comparison — if the underlying counter or flag must be preserved.

\texttt{CALL} saves the address of the instruction \textit{immediately following} the \texttt{CALL} opcode word and its operand (i.e., \texttt{ip} after both words have been consumed) onto the call stack, then sets \texttt{ip} to the target address. \texttt{RET} restores \texttt{ip} from the top of the call stack, resuming execution at the saved address. This mechanism supports up to 128 levels of nested calls before the call stack overflows.

Parameters are conventionally passed on the operand stack, pushed by the caller before \texttt{CALL} and consumed by the function body. Return values are left on the operand stack by the function before \texttt{RET}. This convention is not enforced by the hardware; it is a programming discipline.

An example of a simple function call that doubles its argument:

\begin{lstlisting}
; --- caller ---
PUSH 7       ; word 0: push argument
CALL 6       ; word 2: call function at word 6; return address = word 4
PRINT        ; word 4: print the return value
HALT         ; word 5

; --- function body (starts at word 6) ---
DUP          ; word 6: duplicate argument (a a)
ADD          ; word 7: a + a = 2a
RET          ; word 8: return to word 4
\end{lstlisting}

Output: \texttt{14}

\textbf{Error conditions:} \texttt{JZ} and \texttt{JNZ} halt with \textit{stack underflow} if the stack is empty. \texttt{CALL} halts with \textit{call stack overflow} if 128 frames are already active. \texttt{RET} halts with \textit{call stack underflow} if there is no matching \texttt{CALL} on the call stack. No bounds checking is performed on jump targets; jumping to an out-of-range address produces undefined behaviour.

\subsection{Registers}

The VM provides sixteen general-purpose integer registers, addressed by index 0 through 15. Registers hold their values across instructions and are initialised to zero at VM startup. They offer a way to store and retrieve values without keeping them permanently on the operand stack.

\begin{center}
\begin{tabular}{lll}
\toprule
\textbf{Mnemonic} & \textbf{Stack Effect} & \textbf{Description} \\
\midrule
\texttt{ST <reg>} & \texttt{a --}  & Pop \texttt{a} and store it in register \texttt{reg}. \\
\texttt{LD <reg>} & \texttt{-- a}  & Push the current value of register \texttt{reg}. \\
\bottomrule
\end{tabular}
\end{center}

\texttt{ST} is a destructive write: the value is removed from the operand stack and placed in the register. \texttt{LD} is non-destructive with respect to the register: the register retains its value after the push. A register can be read any number of times with repeated \texttt{LD} instructions.

Registers are particularly useful for loop counters and function arguments that must survive multiple stack operations. Rather than duplicating a value repeatedly across loop iterations, it can be stored once in a register and reloaded as needed:

\begin{lstlisting}
PUSH 10
ST 0       ; register 0 = 10 (loop limit)
PUSH 1
ST 1       ; register 1 = 1  (counter, starts at 1)

; loop body (word 8)
LD 1       ; push counter
PRINT      ; print it
LD 1       ; push counter again
PUSH 1
ADD        ; counter + 1
ST 1       ; store updated counter
LD 1       ; push counter for comparison
LD 0       ; push limit
GT         ; counter > limit?
JNZ 8     ; if not, repeat
HALT
\end{lstlisting}

Output: \texttt{1}, \texttt{2}, \ldots, \texttt{10}

\textbf{Error conditions:} Both \texttt{ST} and \texttt{LD} halt with \textit{invalid register} if the register index is outside the range 0--15. \texttt{ST} additionally halts with \textit{stack underflow} if the stack is empty when the store is attempted.

\subsection{I/O and Control}

\begin{center}
\begin{tabular}{lll}
\toprule
\textbf{Mnemonic} & \textbf{Stack Effect} & \textbf{Description} \\
\midrule
\texttt{PRINT} & \texttt{a --} & Pop the top value and print it to standard output. \\
\texttt{HALT}  & \texttt{--}   & Stop execution; set the running flag to zero. \\
\bottomrule
\end{tabular}
\end{center}

\texttt{PRINT} is the VM's sole output instruction. It pops the top value, formats it as a signed decimal integer, and writes it to \texttt{stdout} followed by a newline. It is intentionally destructive: if the value must be retained, \texttt{DUP} should precede it.

\texttt{HALT} terminates the fetch–decode–execute loop cleanly. Every well-formed program must end with \texttt{HALT}; without it, the VM will continue reading beyond the end of the bytecode array, causing undefined behaviour.

\subsection{Bytecode Encoding Summary}

The table below summarises the word width of each instruction as it appears in the flat bytecode array. Instructions with an operand occupy two words; all others occupy one.

\begin{center}
\begin{tabular}{lcc}
\toprule
\textbf{Instruction(s)} & \textbf{Words} & \textbf{Operand} \\
\midrule
\texttt{ADD, SUB, MULT, DIV, MOD}       & 1 & — \\
\texttt{EQ, LT, GT, AND, OR, NOT}       & 1 & — \\
\texttt{POP, DUP, SWAP, PEEK}           & 1 & — \\
\texttt{PRINT, HALT, RET}               & 1 & — \\
\texttt{PUSH <n>}                       & 2 & Integer literal \\
\texttt{JMP <addr>, JZ <addr>, JNZ <addr>} & 2 & Word index \\
\texttt{CALL <addr>}                    & 2 & Word index \\
\texttt{LD <reg>, ST <reg>}             & 2 & Register index (0--15) \\
\bottomrule
\end{tabular}
\end{center}

Because all values are plain C \texttt{int}s, opcodes and operands share the same representation. The distinction between an opcode word and a data word is purely positional: the VM knows from context (specifically, from the \texttt{operand\_count} field in the assembler's mnemonic table) how many words following an opcode are operands rather than the next opcode.